% =====================================================================
%  Preamble for the LaTeX conversion of
%  Gattringer & Lang, "Quantum Chromodynamics on the Lattice" (LNP 788)
% =====================================================================
\usepackage{amsmath,amssymb,amsthm,mathtools}
\usepackage{graphicx}
\usepackage[svgnames]{xcolor}
\usepackage{microtype}
\usepackage{booktabs}
\usepackage{multirow}
\usepackage{caption}
\usepackage{enumitem}
\usepackage{url}
\usepackage[hidelinks]{hyperref}
\usepackage{cleveref}

% -------- page geometry: reproduce the ~6.1 x 9.25 in lecture-note page --------
\usepackage[papersize={7.0in,9.4in},margin=1.0in,headheight=13pt]{geometry}

% -------- header / footer ---------------------------------------------------
\usepackage{fancyhdr}
\pagestyle{fancy}
\fancyhf{}
\fancyhead[LE]{\small\itshape\leftmark}
\fancyhead[RO]{\small\itshape\rightmark}
\fancyfoot[LE,RO]{\small\thepage}
\renewcommand{\headrulewidth}{0.4pt}
% nicer running heads: "1. The path integral on the lattice" / "1.1 Hilbert space..."
\renewcommand{\chaptermark}[1]{\markboth{\thechapter.\ #1}{}}
\renewcommand{\sectionmark}[1]{\markright{\thesection\ #1}}

% -------- theorem-like environments -----------------------------------------
\newtheorem{definition}{Definition}[chapter]
\newtheorem{remark}[definition]{Remark}
\newtheorem{theorem}[definition]{Theorem}
\newtheorem{lemma}[definition]{Lemma}
\newtheorem{proposition}[definition]{Proposition}
\newtheorem{corollary}[definition]{Corollary}
\newtheorem{example}[definition]{Example}

% -------- physics macros ------------------------------------------------------
\newcommand{\ket}[1]{|#1\rangle}          % |n>
\newcommand{\bra}[1]{\langle#1|}           % <u|
\newcommand{\braket}[2]{\langle#1|#2\rangle}
\newcommand{\bbraket}[3]{\langle#1|#2|#3\rangle}
\newcommand{\vac}{\ket{0}}                 % vacuum state
\newcommand{\tr}{\operatorname{tr}}
\newcommand{\Tr}{\operatorname{Tr}}
\newcommand{\E}{\mathbb{E}}                % expectation value
\newcommand{\dd}{\mathrm{d}}               % upright differential
\newcommand{\D}{\mathrm{d}}
\newcommand{\im}{\mathrm{i}}               % imaginary unit
\renewcommand{\Re}{\operatorname{Re}}
\renewcommand{\Im}{\operatorname{Im}}
\newcommand{\ev}[1]{\left\langle #1 \right\rangle}   % <...> average
\newcommand{\avg}[1]{\left\langle #1 \right\rangle}
\newcommand{\order}[1]{\mathcal{O}\left(#1\right)}
\newcommand{\deriv}[2]{\frac{\dd #1}{\dd #2}}
\newcommand{\pderiv}[2]{\frac{\partial #1}{\partial #2}}
\newcommand{\half}{\tfrac{1}{2}}

% -------- common symbols used throughout the book -----------------------------
\newcommand{\lam}{\Lambda}
\newcommand{\Nc}{N_c}
\newcommand{\Nf}{N_f}
\newcommand{\Ns}{N_s}
\newcommand{\Nt}{N_t}
\newcommand{\mmu}{\mu}
\newcommand{\nnu}{\nu}
\newcommand{\gmu}{\gamma_\mu}

% -------- misc ---------------------------------------------------------------
\newcommand{\sch}{\,\mathrm{Schwinger}}
\setlength{\parskip}{0pt}
\setlength{\parindent}{1.5em}
